% !TEX root = main.tex
% ===========================================================================
% Numbers that cannot (yet) be derived from a committed artefact.
%
% THE RULE: every macro here carries a comment naming where the number came
% from and how it was obtained. If you cannot write that comment, you do not
% know the number well enough to print it.
%
% This file is the exception, not the pattern. Anything derivable belongs in
% scripts/thesis_numbers.py, which regenerates generated/numbers.tex from the
% manifest and the Vivado reports. Whenever a number here becomes derivable -
% a simulation that writes a machine-readable cycle count, an evaluation
% script that writes a JSON - move it there and delete it from here.
%
% Marked [MEASURED] / [DERIVED] / [ESTIMATED] on the project's own standard
% (see DD-019 and \cref{sec:acc-bound}).
% ===========================================================================

% --- verification ---------------------------------------------------------
% Re-measured after DD-036's redeploy, on the leaf-split model (2026-10-22/23):
% the counts did not move, the images did.
% [MEASURED] accel_tb, end-to-end run, 2026-09-22; again on 2026-10-22/23 with
% the redeployed model, in accel_tb, accel_dma_tb and accel_axi_tb. Element
% count of the full-network comparison against the golden vectors. It counts
% elements, not values, so a retrain cannot move it.
\newcommand{\thElements}{6\,895\,780}
% [MEASURED] scripts/gen_program.py output, software/export/program.txt.
\newcommand{\thInstructions}{64}
% [MEASURED] scripts/run_all_sims.sh, all 35 passing on 2026-10-23 with the
% redeployed model. (31 as of 2026-09-24; DD-030 and DD-031 added wgt_dma_tb,
% accel_dma_tb, axi_lite_csr_tb and accel_axi_tb on 2026-09-25.)
\newcommand{\thTestbenches}{35}
% [MEASURED] accel_tb bit-exact on each of the four golden sets, 2026-10-23:
% healthy_9111, blackrot_0346, esca_1744 and blight_4740 - for each class, the
% first test-split photograph in file order that the integer pipeline gets
% right (DD-036). One image per class, so all four argmax outcomes are
% exercised. (Until the redeploy the four were images the leaf-split model
% had been trained or validated on.)
\newcommand{\thImagesVerified}{4}

% --- throughput -----------------------------------------------------------
% [MEASURED] feeders-busy cycle count of one inference: accel_tb, accel_dma_tb
% and accel_axi_tb all print 952,062, on 2026-09-25 and again on 2026-10-22
% with the redeployed model - the schedule does not depend on the weights.
% (The first measurement, 2026-09-13, was 951,833: before the pipelining of
% Chapter~\ref{chap:timing} changed the pipeline and drain depths. This macro
% carried that stale value until 2026-10-22; the figures derived from it round
% the same.)
% TODO: have the testbench write this to a file and move the macro into
% thesis_numbers.py - it is the most-quoted number in the thesis that is still
% transcribed by hand.
\newcommand{\thCycles}{952\,062}
% [MEASURED] compute cycles excluding weight load - the same count as
% \thCycles, kept under its own name where the text sets it against the load.
\newcommand{\thCyclesCompute}{952\,062}
% [MEASURED] accel_dma_tb: whole-network weight load with the real DMA.
\newcommand{\thDmaLoadCycles}{199\,285}
% [MEASURED] the testbench task the DMA replaced, same workload.
\newcommand{\thTaskLoadCycles}{325\,512}
% [DERIVED] 199285 / 952062. NOTE: supersedes the 65% figure in review finding
% C-11, which was computed against a fill rate nobody had measured. 65% still
% stands for features.18.0 alone, the worst single layer.
\newcommand{\thLoadOverheadPct}{17.3}
% [DERIVED] \thCycles / \thClkBoard, i.e. 952062 / 187.5e6 = 5.078 ms.
% NOTE: "Board" means AT THE BOARD'S CLOCK - this is not measured on the
% board, and it is the COMPUTE alone (feeder busy). The measured figure is
% \thBoardMs (the block below, 2026-10-22), which also contains the weight
% loads and the image load. Quote both, and say which is which.
\newcommand{\thLatencyBoard}{5.08}
% [DERIVED] the one-MAC-per-cycle baseline that Section~\ref{sec:parallelism}
% rejects: \thMacs / \thClkBoard, i.e. 299562112 / 187.5e6 = 1.598 s. The
% whole-network total is the right one here, not \thMacsConv - a GAP
% accumulation costs a sequential engine a cycle exactly like a multiply does.
% Quoted in SECONDS. It is 1.6 s, not 1.6 ms; the draft that first carried this
% number was out by 1000x, which made the rejected baseline look three times
% FASTER than \thLatencyBoard and inverted the argument for parallelising.
\newcommand{\thLatencySeq}{1.60}
% [DERIVED] \thMacs / \thCycles, i.e. 299562112 / 952062 = 314.6. A ratio of
% cycle counts, so it is clock-independent; equivalently \thLatencySeq /
% \thLatencyBoard.
\newcommand{\thSpeedupSeq}{315}

% --- measured on the board (DD-035, DD-036) -------------------------------
% [MEASURED] 2026-10-22, docs/lab_notebook.md: `vinecam --image image_q8.bin
% --repeat 100` on the ZCU104 under PetaLinux 2022.2, timed with clock_gettime
% on the A53 from LOAD_IMAGE to the result - the image DMA plus the whole run.
% The only on-board timing there is: the accelerator has no cycle-counter
% register, and the bare-metal app prints no time. Blobs: the export of
% \thCkptSha. The spread is +-3 us, because the schedule does not depend on
% the data.
\newcommand{\thBoardMs}{6.659}
\newcommand{\thBoardMinMs}{6.657}
\newcommand{\thBoardMaxMs}{6.663}
\newcommand{\thBoardRuns}{100}
% [MEASURED] same run: the CPU's copy of one 150,528-byte tensor into the
% uncached DMA region, per image.
\newcommand{\thBoardCopyMs}{0.160}
% [MEASURED] accel_dma_tb, 2026-09-25 (docs/lab_notebook.md): start to done in
% simulation, weight loads included. The image is loaded BEFORE start, so it
% is not in this figure.
\newcommand{\thSimStartDone}{1\,151\,347}
% [DERIVED] \thSimStartDone / \thClkBoard, i.e. 1151347 / 187.5e6 = 6.141 ms.
\newcommand{\thSimStartDoneMs}{6.14}
% [DERIVED] the image load at one pixel per cycle (DD-032): 50,176 / 187.5e6
% = 0.268 ms.
\newcommand{\thImgLoadMs}{0.27}
% [DERIVED] what vinecam times, predicted from simulation: 6.141 + 0.268 =
% 6.408 ms.
\newcommand{\thBoardPredMs}{6.41}
% [DERIVED] measured minus predicted, 6.659 - 6.408 = 0.251 ms, 3.8% of the
% measurement: the PS side (polling, the handshakes around the image load).
% NOT attributed further - say so wherever it is quoted.
\newcommand{\thBoardGapMs}{0.25}
\newcommand{\thBoardGapPct}{3.8}
% [DERIVED] 1000 / \thBoardMs = 150.2 images per second, accelerator alone.
\newcommand{\thBoardIps}{150}
% TODO(BOARD): measured with the previous model's blobs (268f85f714d3). The
% schedule does not depend on the weights, so none of the above should move.
% Re-run --repeat 100 once with the redeployed cards, and say that it did not.
%
% [MEASURED - NOT YET] the live camera loop (lab guide C5, 2026-10-22). It ran
% on the board, but its console status line was not saved. Copy the console
% on the next live run and fill these in. Until then they print as red TODOs.
\newcommand{\thCamFps}{\TODO{live frame rate, from the console}}
\newcommand{\thCamFrontMs}{\TODO{front-end ms per frame, from the console}}

% --- the Tm/Tn design space (DD-013) --------------------------------------
% [DERIVED] scripts/analyze_workload.py, "Tm/Tn sweep at constant P=512".
% The cycle model is HW * ceil(OC/Tm) * ceil(IC/Tn) summed over all
% \thPwLayers{} pointwise layers; utilization is the useful fraction of the
% P=512 MAC slots issued. Re-run the script to regenerate - every figure here
% was reproduced from the current manifest on 2026-10-05 and matches DD-013.
% Suffix names Tm: 32x16 is the decision, the others are the alternatives.
\newcommand{\thPwUtil}{92.1}            % Tm=32 x Tn=16 - DD-013's choice
\newcommand{\thPwUtilSixteen}{88.6}     % 16 x 32
\newcommand{\thPwUtilSixtyFour}{74.9}   % 64 x 8  - the earlier proposal
\newcommand{\thPwUtilEight}{66.1}       % 8  x 64
% [DERIVED] same run, "where the tail-padding goes": 44,758 of 568,088 cycles
% are zero-padded tail tiles. This is 100 - \thPwUtil, from the other side.
\newcommand{\thPwWastePct}{7.9}
% [DERIVED] 2.27 ms pointwise of 4.81 ms modelled total, from the same script's
% depthwise table at Tc=16. Quoted as a SHARE on purpose: the script's absolute
% milliseconds are at 250 MHz, which this design does not reach (\thClkBoard),
% whereas the ratio is clock-independent.
\newcommand{\thPwRuntimePct}{47}
% [ESTIMATED] DD-013, "a pixel-parallel fallback for layers with OC < Tm, which
% recovers only 2.2%". NOTE: unlike every other figure in this block, the sweep
% behind this one is NOT in analyze_workload.py - it exists only in DD-013's
% prose. TODO: add the fallback to the script, or drop the number from the text.
\newcommand{\thPwPixelGain}{2.2}

% --- the depthwise sizing (DD-014) ----------------------------------------
% [DERIVED] scripts/analyze_workload.py, "depthwise cost" table. Shares of the
% modelled inference time, not of the MACs - the whole point of DD-014 is that
% the two differ. Quoted as SHARES for the same reason as \thPwRuntimePct: the
% script's absolute milliseconds are at 250 MHz, a clock this design does not
% reach, while the ratios are clock-independent. Reproduced 2026-10-12.
\newcommand{\thDwShareSeq}{79}          % Tc=1, one output element per cycle
\newcommand{\thDwShareEight}{33}        % Tc=8
\newcommand{\thDwShare}{19}             % Tc=16 - DD-014's choice
\newcommand{\thDwShareThirtyTwo}{11}    % Tc=32, for twice the slices
% [DERIVED] same table: Tc*K*K = 144 multipliers, packed two int8 products per
% DSP48E2. The slice count, not the multiplier count.
\newcommand{\thDwDsp}{72}
% [DERIVED] same script, "stride-2 depthwise layers": features.2/4/7/14.conv.1.0,
% 4 of the \thDwLayers{} depthwise layers, each with an input grid 4x its output.
\newcommand{\thDwStrideTwo}{4}
% [ESTIMATED] DD-014's correction of 2026-09-12: counting OUTPUT pixels instead
% of input pixels understated the depthwise cost by this factor. NOTE: like
% \thPwPixelGain, the before/after comparison is not in analyze_workload.py -
% the script was fixed in place, so only the corrected model survives. The
% factor is recorded in DD-014 alone. TODO: drop it, or keep the old model
% behind a flag so the comparison is reproducible.
\newcommand{\thDwInputFactor}{1.63}

% --- the stem sizing (DD-016) ---------------------------------------------
% [ESTIMATED] DD-016's TS sweep. NOTE: unlike the Tm/Tn and Tc sweeps, this one
% is NOT in analyze_workload.py - the script only computes the stem at one
% element per cycle (sec_depthwise, `stem_c`). The table survives in DD-016 and
% in stem_array.v's header alone.
%
% Quoted as SHARES of the modelled total, which is both clock-independent and
% DD-016's own idiom ("the stem climbs back to 15% of runtime" for TS=4, which
% is 0.55/3.76). Do not quote DD-016's absolute milliseconds: its 401,408
% cycles are printed by the script as 1.61 ms at 250 MHz, but the table calls
% the same figure 1.76 ms. One of the two is wrong and nobody has resolved it.
% TODO: add the TS sweep to analyze_workload.py and settle the discrepancy.
\newcommand{\thStemShareSeq}{35}        % TS=1,  1.76/4.97
\newcommand{\thStemShareFour}{15}       % TS=4,  0.55/3.76
\newcommand{\thStemShare}{10}           % TS=8,  0.35/3.56 - DD-016's choice
\newcommand{\thStemShareSixteen}{7}     % TS=16, 0.25/3.46
% [ESTIMATED] same table, 1.76 / 0.35. It is 5x and NOT 8x: past the knee the
% input stream, not the compute, is the limit - which is the whole of DD-016's
% argument. A draft of \cref{sec:stem-engine} once claimed a linear 8x here.
\newcommand{\thStemSpeedup}{5}
% [DERIVED] TS * CIN * K * K = 8 * 27 = 216 multipliers, packed two int8
% products per DSP48E2.
\newcommand{\thStemDsp}{108}

% --- the requantize-lane count (DD-015) -----------------------------------
% [DERIVED] scripts/analyze_workload.py requant, "requantize parallelism".
% Reproduced 2026-10-12; agrees with DD-015's table. An output tile costs
% max(n_ic, ceil(Tm/R)) cycles, so R trades DSPs against a drain stall. The
% penalties are relative to R=32, which never stalls.
\newcommand{\thRqPenaltyOne}{667}       % R=1,  4,359,316 cycles
\newcommand{\thRqPenaltyEight}{40}      % R=8,    797,016
\newcommand{\thRqPenaltySixteen}{6.6}   % R=16,   605,720 - the knee on cycles
\newcommand{\thRqStalledEight}{7}       % layers stalled at R=8, of \thPwLayers
\newcommand{\thRqStalledSixteen}{1}     % layers stalled at R=16 (features.2.conv.0.0)
% [DERIVED] R=32 lanes x 2 DSP48E2 per 26x32 multiply.
\newcommand{\thRqDsp}{64}
% [DERIVED] the pointwise array's own slices: P = Tm*Tn = 512 int8 MACs packed
% two per DSP48E2. Agrees with DD-015's budget breakdown ("array 256").
\newcommand{\thPwDsp}{256}
%
% [ESTIMATED] out_stage.v's header and docs/lab_notebook.md: the DSP budget
% PREDICTED in September 2026 with and without a shared output stage.
% IMPORTANT: these are predictions from before the pipelining work, not
% measurements. The measured figure for the finished design is
% \thDspAccelAxi{} (\thDspPctAccelAxi\%). Always say "predicted" when quoting
% them, or Chapter~\ref{chap:evaluation} will appear to contradict this one.
\newcommand{\thDspPredUnshared}{634}
\newcommand{\thDspPredShared}{442}

% --- the accumulator width (DD-019) ---------------------------------------
% The DERIVED figures for the deployed model - \thAccBound, \thAccBoundBits,
% \thAccBoundLayer/Channel/Bias, \thAccRange, \thAccMargin,
% \thAccMarginNarrow, \thAccGeom, \thAccUnsafeLayers, \thMaxShift - moved to
% scripts/thesis_numbers.py on 2026-10-22, which replays them from
% software/export/ on every `make numbers`. They used to be copied here by
% hand, and the leaf-split redeploy (DD-036) made every one of them wrong at
% once: DD-019's failure, in prose. What stays here is history, which no
% export can regenerate.
%
% [DERIVED, HISTORY] the previous checkpoint's bound, 268f85f714d3 (the
% DD-019 retrain, deployed until 2026-10-22): the same replay on
% software/export/ as committed at ae0d079 gives 5,094,750 at classifier.1,
% 24 bits, margins 6.6x at ACC_W=26 and 1.6x at 24 - DD-019's own figures -
% and 14 layers over 21 bits. The redeploy moved the bound to a different
% layer and one bit wider; the exporter caught that, nobody had to.
\newcommand{\thAccBoundPrev}{5\,094\,750}
\newcommand{\thAccBoundPrevBits}{24}
\newcommand{\thAccBoundPrevLayer}{classifier.1}
\newcommand{\thAccUnsafeLayersPrev}{14}
% [MEASURED] DD-019: the largest |acc + bias| actually SEEN, one image through
% one checkpoint, at features.3.conv.0.0. Quoted only to be contrasted with the
% bound - it is an observation, and sizing from it is what DD-019 forbids.
\newcommand{\thAccMeasured}{1\,198\,084}
\newcommand{\thAccMeasuredBits}{22}
% NOTE on \thAccGeom (now generated): DD-019's table gives 15,743,028 / 25
% bits for that row, which is the largest CONVOLUTION (~960 taps) and omits the
% classifier (1280 taps). 20,971,520 is the correct figure and fits ACC_W=26.
% TODO: fix DD-019's row.
% [DERIVED, HISTORY] the bias that overflowed a 21-bit accumulator on its own,
% for any input: channel 110 of features.3.conv.0.0 in the DD-019 retrain
% (268f85f714d3) - software/export/features_3_conv_0_0_b.hex as committed at
% ae0d079. The redeploy overwrote that file; this is the figure from git.
\newcommand{\thBadBias}{1\,184\,089}
% [MEASURED] DD-019: elements wrong in the end-to-end golden comparison when
% ACC_W was 21, of \thElements{} total.
\newcommand{\thBadElements}{2\,215\,904}
% [MEASURED] DD-019's controlled pair of synthesis runs, both at the same
% 4.000 ns constraint so ACC_W is the only variable: 32,387 -> 35,531 CLB LUTs.
% DSPs did not move (944 either way), which is the number the decision turned on.
\newcommand{\thAccLutDelta}{3\,144}
\newcommand{\thAccLutPct}{9.7}

% --- the on-chip memories (DD-007, DD-023, DD-029) ------------------------
% UNITS: everything in this block is BINARY (KiB = 1024 B, MiB = 1024 KiB),
% because it sizes memory tiles. Note that \thWeightMB{} above is DECIMAL MB,
% which is how scripts/thesis_numbers.py emits it - so the weight blob is
% quoted here in exact bytes instead, to avoid comparing MB with MiB.
% scripts/gen_program.py itself mixes the two conventions in one print block;
% see the report. Reproduced 2026-10-12 by running gen_program.py.
%
% [DERIVED] gen_program.py, "ACTIVATION POOL PEAK": the static first-fit
% allocator's peak over the whole network, and the largest single tensor.
\newcommand{\thPoolEntries}{50\,176}
\newcommand{\thPoolMiB}{1.53}
\newcommand{\thPoolLargest}{37\,632}
\newcommand{\thPoolLargestMiB}{1.15}
% [DERIVED] the pool is 4 URAM288 tiles wide (ceil(256/72)) by 13 deep
% (ceil(50176/4096)) = 52, which is the whole URAM count the board build
% measures (\thUramBoard{} of 96). The activation pool IS the design's URAM.
\newcommand{\thPoolUramWide}{4}
\newcommand{\thPoolUramDeep}{13}
% [DERIVED] act_buffer.v: a 128-bit pool would cost two cycles per writeback,
% and in features.2.conv.0.0 the store would overtake the compute - this is
% that layer's share of the whole pointwise runtime.
\newcommand{\thNarrowPoolPct}{6.6}
% [DERIVED] act_buffer.v: two fixed ping-pong halves would need twice the
% largest feature map, against the largest resident input+output pair
% (\thPoolMiB{} MiB) that the allocator actually requires.
\newcommand{\thPingPongMiB}{2.3}
%
% [DERIVED] gen_program.py, "weight buffer": worst-case layer occupancy
% (features.18.0) of the buffer's 1024 entries, and the buffer's own capacity
% at TM=32 banks x TN=16 weights x 8 bit = 4096 bit per entry.
\newcommand{\thWgtEntries}{800}
\newcommand{\thWgtDepth}{1024}
\newcommand{\thWgtPct}{78}
\newcommand{\thWgtBufKiB}{512}
\newcommand{\thWgtWorstKiB}{400}
% [DERIVED] the whole weight blob, exact bytes (= \thWeightMB{} decimal MB).
\newcommand{\thWeightBytes}{2\,194\,880}
% [MEASURED] DD-029 / C-11: features.18.0 spends 25,600 fill cycles against
% 39,200 compute, none of it hidden - the WORST SINGLE LAYER. The whole-network
% figure is \thLoadOverheadPct\%, which supersedes this one; quote both only
% with that distinction made, as manual_numbers' \thLoadOverheadPct note says.
\newcommand{\thWgtWorstOverheadPct}{65}
%
% [DERIVED] param_buffer.v: oc-tiles needed against the depth built, at
% TM=32 banks x (BIAS_W 32 + M0_W 32 + SHIFT_W 8) = 72 bit per channel.
\newcommand{\thParamEntries}{40}
\newcommand{\thParamDepth}{64}
% [DERIVED] DD-030: the padded parameter blob, 19,264 channels x 9 B.
\newcommand{\thParamBlobBytes}{173\,376}
%
% [DERIVED] img_buffer.v: 224*224 pixels of CIN=3 bytes, which is both the
% buffer's depth and the DDR image DD-032 defines.
\newcommand{\thImgBytes}{150\,528}
\newcommand{\thImgKiB}{147}

% --- the control unit (T-1, T-2, DD-017) ----------------------------------
% [DERIVED] scripts/gen_program.py's own summary, reproduced 2026-10-15:
%   "instructions: 64 (34 PW, 17 DW, 1 STEM, 10 RES, 1 GAP, 1 LINEAR)"
% The census sums to \thInstructions{}, and top_seq_tb asserts exactly it.
% NOTE on grouping: \thPwLayers{} is 35 because it counts the classifier among
% the pointwise-SHAPED layers; \thOpsPw{} is 34 because LINEAR is its own
% opcode. Both are right; say which grouping is meant.
\newcommand{\thOpsPw}{34}
\newcommand{\thOpsDw}{17}
\newcommand{\thOpsStem}{1}
\newcommand{\thOpsRes}{10}
\newcommand{\thOpsGap}{1}
\newcommand{\thOpsLin}{1}
% [DERIVED] same run: manifest ops that compile to no instruction because they
% move no data - they are liveness annotations for the allocator.
% \thManifestOps{} - \thSaveOps{} = \thInstructions{}.
\newcommand{\thSaveOps}{10}
% [DERIVED] same run, "program memory": \thInstructions{} x 8 words x 4 B.
\newcommand{\thProgBytes}{2\,048}
%
% [DERIVED] hardware/rtl/control/top_seq.v: the FSM's state count, the idle
% cycles left after every op, and the watchdog's width (it trips at 2^20-1
% cycles, shared across S_LOAD / S_WAIT / S_RUN).
\newcommand{\thSeqStates}{8}
% [DERIVED] accel_top.v:147 and accel_axi.v:106 both set GAP_CYC = OS_LAT, so
% the built design idles \thOsLat{} cycles, not the 4 that top_seq.v carries as
% a MODULE DEFAULT. Same trap as act_buffer's DEPTH: the default is not what is
% instantiated. T-10 made it derived; top_seq.v's own header still says
% "64*4 cycles over the whole network", which is stale. TODO: fix that header.
\newcommand{\thGapCyc}{5}
\newcommand{\thWdogBits}{20}
% [DERIVED] pw_feeder.v: DRAIN_CYC = ARR_LAT + OS_LAT = 7 + 5. This is how far
% `busy` must outlive the address generator on the pointwise path.
% NOTE: DD-017 says "~7 extra cycles", which predates the L-1/L-3 pipelining
% that deepened both terms. 12 is the current structural figure; do not quote
% the 7. TODO: correct DD-017.
\newcommand{\thDrainPw}{12}

% --- the latency contract (docs/latency_map.md) ---------------------------
% [DERIVED] the six constants every other depth in the design is sized from,
% each a localparam in the module named, each reproduced from the RTL on
% 2026-10-20. The "before" column is the pre-pipelining value recorded in
% docs/latency_map.md section 1; L-1 pipelined the MAC trees (phases 2a-2c),
% L-3 the requantizer (2d), A-1 added the array input register (phase 3).
\newcommand{\thMacLat}{6}        \newcommand{\thMacLatOld}{1}
\newcommand{\thArrLat}{7}        % = 1 + MAC_LAT, after A-1
\newcommand{\thDwLat}{5}         \newcommand{\thDwLatOld}{0}
\newcommand{\thStLat}{6}         \newcommand{\thStLatOld}{0}
\newcommand{\thRqLat}{4}         \newcommand{\thRqLatOld}{1}
\newcommand{\thOsLat}{5}         \newcommand{\thOsLatOld}{2}
% [DERIVED] the five drain counters, which are what those constants are FOR.
% Each is now a localparam derived from the stages it covers; the "old" values
% are the hand-set literals docs/latency_map.md section 3 records, which
% carried 3-8 cycles of unexplained margin. \thDrainPw{} is defined above.
\newcommand{\thDrainPwOld}{7}
\newcommand{\thDrainDw}{11}      \newcommand{\thDrainDwOld}{7}
\newcommand{\thDrainStem}{12}    \newcommand{\thDrainStemOld}{15}
\newcommand{\thDrainRes}{5}      \newcommand{\thDrainResOld}{6}
\newcommand{\thDrainGap}{6}      \newcommand{\thDrainGapOld}{7}
% [DERIVED] docs/latency_map.md: the watchdog is sized against the longest
% legitimate operation (the stem) rather than against any latency.
\newcommand{\thStemCycles}{100\,801}

% --- per-module cost ------------------------------------------------------
% [MEASURED] out-of-context synthesis, DD-031. Neither uses a DSP or a BRAM.
\newcommand{\thLutCsr}{173}
\newcommand{\thFfCsr}{354}
\newcommand{\thLutDma}{1\,720}
\newcommand{\thFfDma}{1\,396}

% --- accuracy -------------------------------------------------------------
% Moved to scripts/thesis_numbers.py (\thAccFloat, \thAccInt, \thAccDrop and
% the rest), which reads the report evaluate.py writes for the held-out test
% split of data/grape_leafsplit - DD-034. Neither 100.00% figure in the
% repository's history is a result: both come from data/grape, which leaked
% rotated copies of each image across its split. data/grape_clean, the first
% fix, still leaked the other three photos of each leaf.
